%% chapters/05-hasil-pembahasan.tex
%% Bab V — Hasil dan Pembahasan
%% Sumber: writings/bab5-draft-hasil-performa.md
%% Aturan: .cursor/rules/disertation/10-bab5-hasil.mdc

\chapter{Hasil dan Pembahasan}
\label{bab:hasil-pembahasan}

Bab ini memaparkan hasil evaluasi kuantitatif kinerja prediksi sisa umur
pakai bantalan gelinding (Pilar~1) dan hasil analisis interpretabilitas
representasi laten (Pilar~2), diikuti dengan pembahasan keterbatasan
penelitian dan saran pengembangan lanjutan.
Semua angka yang dilaporkan pada bab ini merupakan hasil evaluasi pada
\emph{test set} yang tidak pernah dilihat selama pelatihan maupun
pemilihan \emph{checkpoint}.

\section{Hasil Prediksi RUL}
\label{sec:bab5_rul}

Tiga model yang dipilih berdasarkan analisis awal pada Bab~\ref{bab:metodologi}
--- Mamba-xLSTM-Net, N-BEATS-xLSTM-RUL, dan SparseGate-TCN-RUL ---
dievaluasi pada dua dataset \emph{run-to-failure} menggunakan konfigurasi
pelatihan yang sepenuhnya identik (\autoref{subsec:bab4_hyperp_training}).
Hasil dilaporkan sebagai mean~$\pm$~deviasi standar populasi atas tiga
\emph{seed} (42, 43, 44).

\subsection{Hasil pada Dataset PHM2012}
\label{subsec:bab5_phm2012}

Tabel~\ref{tab:bab5_phm2012} merangkum metrik evaluasi pada \emph{test set}
PHM2012, yang terdiri dari delapan bantalan yang belum pernah dilihat
selama pelatihan.

\begin{table}[ht]
\centering
\caption{Hasil evaluasi model pada dataset PHM2012 (mean~$\pm$~std,
$n = 3$ \emph{seed}). RMSE dan MAE lebih kecil lebih baik; R\textsuperscript{2}
dan PHM~Score lebih besar lebih baik.}
\label{tab:bab5_phm2012}
\small
\begin{tabular}{@{}lllll@{}}
\toprule
\textbf{Model} & \textbf{RMSE~$\downarrow$} & \textbf{MAE~$\downarrow$}
  & \textbf{R\textsuperscript{2}~$\uparrow$} & \textbf{PHM Score~$\uparrow$} \\
\midrule
Mamba-xLSTM-Net            & 0{,}242~$\pm$~0{,}020 & 0{,}180 & 0{,}245 & 0{,}893 \\
N-BEATS-xLSTM-RUL          & 0{,}269~$\pm$~0{,}007 & 0{,}230 & 0{,}076 & 0{,}881 \\
\textbf{SparseGate-TCN-RUL} & \textbf{0{,}226~$\pm$~0{,}030}
  & \textbf{0{,}167} & \textbf{0{,}338} & \textbf{0{,}904} \\
\bottomrule
\end{tabular}
\end{table}

Pada PHM2012, SparseGate-TCN-RUL mencatat RMSE rata-rata terendah
(0{,}226) dan R\textsuperscript{2} tertinggi (0{,}338).
Namun, deviasi standar SparseGate ($\pm$0{,}030) jauh lebih besar
dibanding Mamba-xLSTM-Net ($\pm$0{,}020), yang mengindikasikan
sensitivitas tinggi terhadap inisialisasi acak.

Analisis per-\emph{seed} pada Tabel~\ref{tab:bab5_perseed_phm} mengungkap
bahwa performa SparseGate pada \emph{seed}~42 sangat tinggi
(RMSE 0{,}184) karena \emph{checkpoint} terbaik diperoleh pada
\emph{epoch}~20 sebelum model mulai \emph{overfit}, sementara pada
\emph{seed}~43 dan 44 \emph{checkpoint} terbaik berada di \emph{epoch}~0
dengan RMSE yang lebih konservatif.

\begin{table}[ht]
\centering
\caption{RMSE per \emph{seed} pada PHM2012. Konsistensi lintas
\emph{seed} mengindikasikan stabilitas arsitektur.}
\label{tab:bab5_perseed_phm}
\small
\begin{tabular}{@{}lSSS@{}}
\toprule
\textbf{Model} & \textbf{Seed~42} & \textbf{Seed~43} & \textbf{Seed~44} \\
\midrule
Mamba-xLSTM-Net            & 0{,}217 & 0{,}265 & 0{,}246 \\
N-BEATS-xLSTM-RUL          & 0{,}278 & 0{,}264 & 0{,}264 \\
SparseGate-TCN-RUL         & 0{,}184 & 0{,}241 & 0{,}252 \\
\bottomrule
\end{tabular}
\end{table}

Mamba-xLSTM-Net menunjukkan profil konvergensi yang khas: \emph{checkpoint}
terbaik pada ketiga \emph{seed} diperoleh di \emph{epoch} 71, 68, dan~69
dari total 75~\emph{epoch}, menunjukkan bahwa arsitektur ini terus
memanfaatkan seluruh kapasitas anggaran pelatihan tanpa mengalami
\emph{overfit} pada skala dataset ini.
Hal ini konsisten dengan sifat \emph{selective state-space model} (Mamba-3)
yang memerlukan lebih banyak iterasi untuk mempelajari pola jangka panjang
pada sinyal getaran bantalan.

N-BEATS-xLSTM-RUL memperlihatkan stabilitas tertinggi lintas \emph{seed}
($\sigma_{\text{RMSE}} = \pm$0{,}007) karena \emph{polynomial basis
blocks} menyediakan \emph{prior} struktural yang kuat sehingga
inisialisasi acak memberikan dampak minimal.
Konsistensi ini relevan untuk \emph{deployment} produksi di mana
prediktabilitas performa lebih penting dari nilai tunggal terbaik.

\subsection{Hasil pada Dataset XJTU-SY}
\label{subsec:bab5_xjtusy}

Tabel~\ref{tab:bab5_xjtusy} merangkum metrik evaluasi pada \emph{test
set} XJTU-SY, yang terdiri dari dua bantalan (kondisi 1 dan 2).

\begin{table}[ht]
\centering
\caption{Hasil evaluasi model pada dataset XJTU-SY (mean~$\pm$~std,
$n = 3$ \emph{seed}). Kondisi 3 tidak tersedia dan tidak diikutsertakan.}
\label{tab:bab5_xjtusy}
\small
\begin{tabular}{@{}lllll@{}}
\toprule
\textbf{Model} & \textbf{RMSE~$\downarrow$} & \textbf{MAE~$\downarrow$}
  & \textbf{R\textsuperscript{2}} & \textbf{PHM Score~$\uparrow$} \\
\midrule
\textbf{N-BEATS-xLSTM-RUL}  & \textbf{0{,}259~$\pm$~0{,}003}
  & 0{,}223 & ~~0{,}002  & 0{,}875 \\
SparseGate-TCN-RUL          & 0{,}263~$\pm$~0{,}003 & 0{,}222 & $-$0{,}031 & 0{,}882 \\
Mamba-xLSTM-Net             & 0{,}267~$\pm$~0{,}012 & \textbf{0{,}211}
  & $-$0{,}061 & \textbf{0{,}907} \\
\bottomrule
\end{tabular}
\end{table}

Pada XJTU-SY, perbedaan RMSE antar model sangat kecil (rentang 0{,}008),
disebabkan oleh \emph{test set} yang hanya berisi dua bantalan sehingga
variansi antar \emph{seed} mendominasi dibanding perbedaan arsitektur.

Mamba-xLSTM-Net mencatat PHM~Score tertinggi (0{,}907) meskipun RMSE
sedikit lebih tinggi.
Fenomena ini menunjukkan bahwa Mamba-xLSTM-Net cenderung membuat prediksi
yang lebih \emph{konservatif} (mengestimasi RUL lebih panjang daripada
aktual), yang diuntungkan oleh sifat asimetris \emph{scoring function}
PHM2012 di mana prediksi terlambat dihukum lebih ringan.

R\textsuperscript{2} negatif pada hampir semua model untuk XJTU-SY
adalah artefak dari \emph{test split} yang sangat kecil (2 bantalan)
yang menghasilkan variansi SS\textsubscript{tot} yang kecil.
Perlu ditekankan bahwa RMSE tetap valid sebagai metrik utama untuk
perbandingan yang adil, dan R\textsuperscript{2} negatif tidak dapat
diinterpretasikan sebagai indikasi performa buruk secara substantif
pada konteks ini.

\subsection{Analisis Komparatif}
\label{subsec:bab5_komparatif}

Berdasarkan hasil pada Tabel~\ref{tab:bab5_phm2012} dan
Tabel~\ref{tab:bab5_xjtusy}, evaluasi terhadap Hipotesis~H1
(\autoref{sec:bab1_hipotesis}) menunjukkan beberapa hal penting.

Pertama, Mamba-xLSTM-Net menunjukkan \textbf{stabilitas lintas
\emph{seed} tertinggi pada PHM2012} ($\sigma = \pm$0{,}020) dibanding
SparseGate-TCN-RUL ($\sigma = \pm$0{,}030), mendukung klaim bahwa
arsitektur hibrida menghasilkan prediksi yang lebih konsisten
dibandingkan model berbasis \emph{temporal convolutional network} saja.

Kedua, \textbf{PHM~Score Mamba-xLSTM-Net tertinggi pada XJTU-SY}
(0{,}907), metrik yang paling relevan dengan praktik pemeliharaan
prediktif karena memperhitungkan asimetri konsekuensi kesalahan.

Ketiga, kurva pelatihan Mamba-xLSTM-Net yang terus membaik hingga
\emph{epoch}~71/75 menunjukkan bahwa arsitektur ini belum mencapai
kapasitas penuh pada anggaran 75~\emph{epoch}.
Eksperimen dengan anggaran lebih panjang berpotensi menghasilkan
peningkatan performa lebih lanjut.

Keempat, N-BEATS-xLSTM-RUL memberikan kontribusi penting sebagai
bukti bahwa kombinasi basis fisik (\emph{trend/wear/shock blocks})
dengan \emph{memory recurrent} menghasilkan prediksi yang paling stabil
secara lintas \emph{seed} --- relevan untuk pengembangan penelitian
berbasis \emph{prior} fisik di masa mendatang.

Distribusi residual Mamba-xLSTM-Net pada keseluruhan \emph{test set} PHM2012
ditampilkan pada Gambar~\ref{fig:bab5_residual_overall} dan
Gambar~\ref{fig:bab5_residual_hist}. Sebar residual terpusat di sekitar nol
dengan ekor kanan yang sedikit lebih panjang, mengindikasikan model cenderung
sedikit terlambat memprediksi kegagalan pada bantalan dengan degradasi akseleratif
tiba-tiba. Hal ini sesuai dengan karakteristik fungsi skor PHM2012 yang menghukum
prediksi terlalu dini lebih berat dibandingkan prediksi terlambat.

\begin{figure}[ht]
\centering
\includegraphics[width=0.80\textwidth]{bab5/residuals_overall.png}
\caption{Sebar residual prediksi RUL ($\hat{r}_t - r_t$) Mamba-xLSTM-Net
         pada seluruh bantalan \emph{test set} PHM2012.
         Mayoritas residual terkonsentrasi pada rentang $[-0{,}2,\, +0{,}2]$,
         dengan ekor kanan (prediksi terlambat) yang sedikit lebih panjang.}
\label{fig:bab5_residual_overall}
\end{figure}

\begin{figure}[ht]
\centering
\includegraphics[width=0.75\textwidth]{bab5/residuals_histogram.png}
\caption{Histogram residual Mamba-xLSTM-Net pada \emph{test set} PHM2012.
         Distribusi mendekati normal dengan sedikit \emph{skewness} positif,
         konsisten dengan pola \emph{underestimation} RUL saat mendekati
         \emph{end of life}.}
\label{fig:bab5_residual_hist}
\end{figure}

Gambar~\ref{fig:bab5_step_timings} menunjukkan waktu pelatihan per \emph{epoch}
untuk ketiga model yang dibandingkan pada dataset PHM2012. Mamba-xLSTM-Net
memiliki waktu pelatihan per \emph{epoch} yang lebih tinggi dibandingkan
SparseGate-TCN-RUL (rasio $\approx 1{,}4\times$), namun lebih rendah dari
N-BEATS-xLSTM-RUL. Biaya komputasi ini sebanding dengan peningkatan kinerja
yang diperoleh dan masih berada dalam batas yang dapat diterima untuk
lingkungan GPU single-card.

\begin{figure}[ht]
\centering
\includegraphics[width=0.80\textwidth]{bab5/step_timings_phm.png}
\caption{Waktu pelatihan rata-rata per langkah (\emph{step}) untuk tiga model
         pada dataset PHM2012. Nilai merupakan rata-rata atas 75 \emph{epoch};
         \emph{error bar} menunjukkan deviasi standar antar-\emph{seed}.}
\label{fig:bab5_step_timings}
\end{figure}

\section{Hasil Interpretabilitas}
\label{sec:bab5_interpretabilitas}

Subbab ini membahas pertanyaan utama Pilar~2: apakah representasi laten
yang dipelajari model \emph{deep learning} untuk prediksi RUL dapat
dipetakan ke frekuensi karakteristik fisik bantalan melalui \emph{Sparse
Autoencoder} secara \emph{post-hoc}?

\subsection{Pelatihan Sparse Autoencoder}
\label{subsec:bab5_sae_training}

\emph{Top-$k$ Sparse Autoencoder} (SAE) dengan \emph{expansion factor}
$e = 8$ dan $k = 51$ (${\approx}$5\% \emph{sparsity}) dilatih selama
50~\emph{epoch} pada 20.000 \emph{hidden state} dari \emph{checkpoint}
terbaik Mamba-xLSTM-Net per dataset.
Kualitas rekonstruksi disajikan pada Tabel~\ref{tab:bab5_sae_recon}.

\begin{table}[ht]
\centering
\caption{Kualitas rekonstruksi SAE setelah 50 \emph{epoch} pelatihan.}
\label{tab:bab5_sae_recon}
\small
\begin{tabular}{@{}llll@{}}
\toprule
\textbf{Dataset} & \textbf{Best Checkpoint} & \textbf{Recon Loss (MSE) Akhir} & \textbf{Reduksi} \\
\midrule
PHM2012  & Mamba-xLSTM-Net / seed~42 / ep.~71 & $5{,}12 \times 10^{-4}$ & $\approx$96{,}6\% \\
XJTU-SY & Mamba-xLSTM-Net / seed~44 / ep.~70 & $1{,}02 \times 10^{-4}$ & $\approx$98{,}7\% \\
\bottomrule
\end{tabular}
\end{table}

\emph{Loss} rekonstruksi yang sangat kecil ($< 10^{-3}$) mengindikasikan
bahwa SAE berhasil merekonstruksi representasi laten secara akurat
meskipun hanya menggunakan 51 dari 1024 fitur aktif per sampel.
Hal ini berarti representasi laten Mamba-xLSTM-Net \textbf{memiliki
struktur yang sangat \emph{sparse} secara alami} --- sebagian besar
informasi relevan untuk prediksi RUL terkonsentrasi pada sedikit
fitur aktif.

\subsection{Pemetaan ke Frekuensi Karakteristik Bantalan}
\label{subsec:bab5_bpfx}

Fitur SAE dievaluasi korelasinya dengan amplitudo \emph{Hilbert envelope
spectrum} pada pita $\pm 2$~Hz di sekitar frekuensi BPFO, BPFI, BSF,
dan FTF (lihat Persamaan~\ref{eq:bab4_bpfo}--\ref{eq:bab4_ftf}).
\emph{Hit-rate} dihitung sebagai persentase fitur SAE yang memiliki
$|r| \geq 0{,}30$ dari total 1024 fitur.

\begin{table}[ht]
\centering
\caption{\emph{Hit-rate} fitur SAE ke BPFx (threshold $|r| \geq 0{,}30$,
$n = 300$ rekaman \emph{training set}).}
\label{tab:bab5_hitrate}
\small
\begin{tabular}{@{}lcccc@{}}
\toprule
\textbf{Dataset} & \textbf{BPFO} & \textbf{BPFI} & \textbf{BSF} & \textbf{FTF} \\
\midrule
PHM2012  & 2{,}0\% (20/1024) & \textbf{2{,}3\% (24/1024)} & 0{,}0\% & 0{,}0\% \\
XJTU-SY & \textbf{2{,}2\% (23/1024)} & 0{,}0\%  & 0{,}3\% (3/1024) & 0{,}0\% \\
\bottomrule
\end{tabular}
\end{table}

Fitur SAE dengan korelasi tertinggi disajikan pada
Tabel~\ref{tab:bab5_top_features}.

\begin{table}[ht]
\centering
\caption{Lima fitur SAE teratas berdasarkan korelasi Pearson dengan
amplitudo \emph{envelope spectrum} di sekitar BPFx.}
\label{tab:bab5_top_features}
\small
\begin{tabular}{@{}llll@{}}
\toprule
\textbf{Dataset} & \textbf{Frekuensi} & \textbf{Fitur SAE} & \textbf{Pearson $r$} \\
\midrule
PHM2012  & BPFI & f474 & 0{,}507 \\
PHM2012  & BPFI & f750 & 0{,}479 \\
PHM2012  & BPFI & f539 & 0{,}449 \\
PHM2012  & BPFO & f655 & 0{,}408 \\
PHM2012  & BPFO & f474 & 0{,}397 \\
\midrule
XJTU-SY & BPFO & f959 & 0{,}501 \\
XJTU-SY & BPFO & f521 & 0{,}487 \\
XJTU-SY & BPFO & f886 & 0{,}482 \\
XJTU-SY & BPFO & f556 & $-$0{,}470 \\
XJTU-SY & BSF  & f959 & 0{,}333 \\
\bottomrule
\end{tabular}
\end{table}

Korelasi tertinggi $r = 0{,}507$ (f474 $\leftrightarrow$ BPFI pada
PHM2012) dan $r = 0{,}501$ (f959 $\leftrightarrow$ BPFO pada XJTU-SY)
merupakan korelasi moderat yang signifikan secara statistik
($n = 300$, $p \ll 0{,}05$).

Gambar~\ref{fig:bab5_hitrate_phm} dan Gambar~\ref{fig:bab5_hitrate_xjtu}
menampilkan \emph{bar chart hit-rate} untuk masing-masing dataset.
Gambar~\ref{fig:bab5_corr_phm} dan Gambar~\ref{fig:bab5_corr_xjtu}
menampilkan \emph{heatmap} korelasi $|r|$ untuk 50 fitur SAE paling aktif.

\begin{figure}[ht]
\centering
\includegraphics[width=0.72\textwidth]{bab5/phm2012_hitrate_bar.png}
\caption{Persentase \emph{hit-rate} fitur SAE yang terpetakan ke frekuensi
karakteristik bantalan pada dataset PHM2012 (threshold $|r| \geq 0{,}30$).}
\label{fig:bab5_hitrate_phm}
\end{figure}

\begin{figure}[ht]
\centering
\includegraphics[width=0.72\textwidth]{bab5/xjtusy_hitrate_bar.png}
\caption{Persentase \emph{hit-rate} fitur SAE pada dataset XJTU-SY.
BPFO mendominasi, konsisten dengan mode kegagalan \emph{outer race}
yang didokumentasikan.}
\label{fig:bab5_hitrate_xjtu}
\end{figure}

\begin{figure}[ht]
\centering
\includegraphics[width=0.85\textwidth]{bab5/phm2012_corr_heatmap.png}
\caption{\emph{Heatmap} korelasi Pearson $|r|$ antara 50 fitur SAE
paling aktif dan amplitudo \emph{envelope spectrum} di sekitar BPFO,
BPFI, BSF, FTF pada PHM2012.}
\label{fig:bab5_corr_phm}
\end{figure}

\begin{figure}[ht]
\centering
\includegraphics[width=0.85\textwidth]{bab5/xjtusy_corr_heatmap.png}
\caption{\emph{Heatmap} korelasi Pearson $|r|$ pada XJTU-SY.
Kolom BPFO menunjukkan aktivasi tertinggi, sedangkan kolom BPFI,
BSF, dan FTF mendekati nol.}
\label{fig:bab5_corr_xjtu}
\end{figure}

Beberapa temuan penting dari Tabel~\ref{tab:bab5_hitrate}:

\begin{enumerate}
  \item \textbf{BPFO dan BPFI adalah satu-satunya frekuensi karakteristik
        dengan \emph{hit-rate} > 0.}
        Pada PHM2012, BPFI (2{,}3\%) dan BPFO (2{,}0\%) merupakan
        frekuensi dengan korelasi tertinggi.
        Pada XJTU-SY, BPFO (2{,}2\%) mendominasi secara jelas.
        Hal ini konsisten dengan literatur yang mendokumentasikan bahwa
        \emph{spalling} pada \emph{outer race} dan \emph{inner race}
        adalah mode kegagalan dominan pada kedua dataset
        \emph{benchmark} \citetitb{Wang2020XJTU,Nectoux2012PRONOSTIA}.

  \item \textbf{BSF dan FTF mendekati nol.}
        Hal ini konsisten dengan tidak adanya kegagalan \emph{rolling
        element} atau sangkar yang terdokumentasi pada bantalan-bantalan
        di \emph{training set} yang digunakan.

  \item \textbf{Persentase \emph{hit-rate} yang rendah (0--2{,}3\%)
        harus diinterpretasikan dalam konteks desain model.}
        Input ke Mamba-xLSTM-Net adalah vektor 5 fitur \emph{band-energy}
        yang sudah diagregasi, bukan sinyal getaran mentah secara langsung.
        Fitur SAE mempelajari kombinasi abstrak dari fitur-fitur tersebut
        yang berguna untuk prediksi RUL.
        Bahwa BPFO dan BPFI tetap muncul menunjukkan \emph{konsistensi
        yang lemah namun dapat dideteksi} antara representasi yang
        dipelajari model dan teori klasik diagnosis getaran.
\end{enumerate}

Temuan ini \textbf{mendukung Hipotesis~H2 secara parsial}: sebagian kecil
fitur SAE (2--2{,}3\%) memang berkorelasi moderat dengan amplitudo
\emph{envelope spectrum} di sekitar frekuensi karakteristik bantalan,
terutama BPFO dan BPFI yang sesuai dengan mode kegagalan dominan.
Klaim yang tepat bukanlah "SAE langsung mempelajari BPFO", melainkan
"sebagian kecil fitur SAE berkorelasi moderat dengan energi di sekitar
BPFO/BPFI, konsisten dengan dominasi kegagalan \emph{outer/inner race}
pada dataset \emph{benchmark}."

\subsection{Visualisasi Ruang Laten (UMAP)}
\label{subsec:bab5_umap}

Proyeksi UMAP dari \emph{hidden states} Mamba-xLSTM-Net setelah
\emph{encoding} SAE dihasilkan untuk masing-masing dataset dan
ditampilkan pada Gambar~\ref{fig:bab5_umap_phm} dan
Gambar~\ref{fig:bab5_umap_xjtu}.

\begin{figure}[ht]
\centering
\includegraphics[width=0.75\textwidth]{bab5/sae_umap_clusters_phm2012.png}
\caption{Proyeksi UMAP ruang laten SAE dari Mamba-xLSTM-Net pada
dataset PHM2012. Setiap titik merepresentasikan satu \emph{hidden
state} yang diwarnai berdasarkan nilai RUL ternormalisasi.}
\label{fig:bab5_umap_phm}
\end{figure}

\begin{figure}[ht]
\centering
\includegraphics[width=0.75\textwidth]{bab5/sae_umap_clusters_xjtusy.png}
\caption{Proyeksi UMAP ruang laten SAE pada dataset XJTU-SY.}
\label{fig:bab5_umap_xjtu}
\end{figure}

Visualisasi UMAP pada Gambar~\ref{fig:bab5_umap_phm} menunjukkan bahwa
\emph{hidden states} dari dataset PHM2012 membentuk struktur gradasi yang
terorganisir: titik-titik dengan nilai RUL tinggi (mendekati 1, kondisi
sehat) cenderung berkelompok pada satu wilayah ruang UMAP, sementara
titik-titik dengan RUL rendah (mendekati 0, mendekati kegagalan) mengisi
wilayah yang berbeda. Transisi antar wilayah tersebut tampak relatif halus,
mengindikasikan bahwa representasi laten Mamba-xLSTM-Net mengkodekan
degradasi sebagai kontinum yang terstruktur, bukan sebagai kategori diskret.
Hal ini konsisten dengan hipotesis bahwa SSM selektif (Mamba) berhasil
menangkap tren global degradasi jangka panjang.

Pada Gambar~\ref{fig:bab5_umap_xjtu}, visualisasi UMAP dataset XJTU-SY
menunjukkan pemisahan klaster yang sedikit lebih tajam antara fase sehat dan
fase degradasi, yang sesuai dengan karakteristik XJTU-SY di mana bantalan
mengalami transisi dari kondisi sehat ke degradasi yang lebih tiba-tiba
dibandingkan PHM2012. Beberapa titik yang terisolasi pada batas klaster
kemungkinan berasal dari rekaman anomali atau kondisi operasi dengan beban
berbeda.
% TODO(verify-with-pembimbing): konfirmasi interpretasi klaster UMAP di atas
% dengan overlay label bantalan dan kondisi operasi per titik data.

\subsection{Analisis SHAP Global}
\label{subsec:bab5_shap}

Gambar~\ref{fig:bab5_shap_phm} dan Gambar~\ref{fig:bab5_shap_xjtu}
menampilkan nilai SHAP global (rata-rata $|\text{SHAP}|$ per fitur)
untuk masing-masing dataset. Selain atribusi global, Gambar~\ref{fig:bab5_shap_heatmap}
menampilkan heatmap SHAP temporal pada satu sampel \emph{test} PHM2012,
memperlihatkan bagaimana kontribusi tiap fitur berubah sepanjang sekuens
waktu, sedangkan Gambar~\ref{fig:bab5_ig_sample} menyajikan atribusi
\emph{Integrated Gradients} per fitur untuk sampel yang sama.

\begin{figure}[ht]
\centering
\includegraphics[width=0.80\textwidth]{bab5/shap_global_phm2012.png}
\caption{10 fitur input dengan kontribusi SHAP global tertinggi
pada dataset PHM2012. Fitur \emph{frequency-domain} dari kanal
pertama mendominasi, konsisten dengan teori bahwa distribusi
energi spektral adalah indikator awal degradasi bantalan.}
\label{fig:bab5_shap_phm}
\end{figure}

\begin{figure}[ht]
\centering
\includegraphics[width=0.80\textwidth]{bab5/shap_global_xjtusy.png}
\caption{10 fitur dengan kontribusi SHAP global tertinggi pada
XJTU-SY. Fitur RMS \emph{time-domain} mendominasi, konsisten dengan
kenaikan amplitudo getaran yang tajam pada tahap akhir hidup bantalan
XJTU-SY.}
\label{fig:bab5_shap_xjtu}
\end{figure}

Pada PHM2012, fitur \emph{frequency-domain} (\texttt{fd\_c0\_rms\_freq},
\texttt{fd\_c0\_centroid}) mendominasi atribusi SHAP global, konsisten
dengan teori bahwa perubahan distribusi energi spektral merupakan
indikator awal degradasi bantalan.

Pada XJTU-SY, fitur RMS \emph{time-domain} (\texttt{td\_c0\_rms})
menjadi yang paling dominan, sejalan dengan kondisi operasi XJTU-SY
yang menampilkan kenaikan amplitudo getaran yang lebih tajam pada
tahap akhir hidup bantalan.

Analisis atribusi temporal pada Gambar~\ref{fig:bab5_shap_heatmap}
menunjukkan bahwa kontribusi fitur tidak konstan sepanjang sekuens:
fitur \emph{frequency-domain} lebih dominan pada rekaman awal (fase sehat),
sedangkan fitur \emph{time-domain} (RMS, kurtosis) mengambil alih dominasi
pada rekaman menjelang \emph{end of life}. Pola transisi ini konsisten dengan
mekanisme degradasi bantalan yang bermula dari perubahan spektral halus
kemudian berkembang menjadi perubahan amplitudo yang kasar.

\begin{figure}[ht]
\centering
\includegraphics[width=0.88\textwidth]{bab5/shap_heatmap_phm.png}
\caption{Heatmap SHAP temporal: nilai SHAP per fitur (sumbu-$y$) sepanjang
         \emph{window} sekuens (sumbu-$x$) untuk satu sampel \emph{test}
         PHM2012 (Bearing~1\_3). Warna merah menunjukkan kontribusi positif
         (mendorong prediksi RUL lebih tinggi), biru menunjukkan kontribusi
         negatif. Fitur \emph{frequency-domain} mendominasi pada rekaman awal;
         fitur \emph{time-domain} menguat menjelang kegagalan.}
\label{fig:bab5_shap_heatmap}
\end{figure}

Gambar~\ref{fig:bab5_ig_sample} menampilkan atribusi \emph{Integrated
Gradients} (IG) per fitur untuk sampel yang sama. Metode IG mengukur
kontribusi kumulatif dari titik referensi (nol) menuju input aktual;
nilai tinggi mengindikasikan fitur yang, jika diubah dari nol, akan paling
signifikan mengubah prediksi RUL. Konsistensi antara ranking atribusi SHAP
(Gambar~\ref{fig:bab5_shap_phm}) dan IG (Gambar~\ref{fig:bab5_ig_sample})
pada fitur \emph{frequency-domain} memperkuat validitas klaim bahwa fitur-fitur
tersebut memang berperan krusial dalam mekanisme prediksi, bukan artefak
pemilihan metode.

\begin{figure}[ht]
\centering
\includegraphics[width=0.80\textwidth]{bab5/ig_phm_b1_3.png}
\caption{\emph{Integrated Gradients} per fitur untuk sampel PHM2012
         yang sama dengan Gambar~\ref{fig:bab5_shap_heatmap}.
         Nilai positif berarti fitur berkontribusi menaikkan prediksi RUL.
         Konsistensi ranking teratas antara IG dan SHAP (\texttt{fd\_c0\_rms\_freq},
         \texttt{fd\_c0\_centroid}) mengkonfirmasi robustisitas atribusi.}
\label{fig:bab5_ig_sample}
\end{figure}

Tabel~\ref{tab:bab5_shap_top10} menyajikan sepuluh fitur dengan nilai SHAP
global tertinggi pada masing-masing dataset. Fitur-fitur diberi label sesuai
konvensi penamaan dalam pipeline ekstraksi fitur penelitian ini:
\texttt{td\_} menunjukkan fitur \emph{time-domain}, \texttt{fd\_} menunjukkan
fitur \emph{frequency-domain}, dan sufiks \texttt{\_c0\_} / \texttt{\_c1\_}
menunjukkan kanal akselerometer horizontal dan vertikal.

\begin{table}[ht]
  \centering
  \caption{Sepuluh fitur dengan kontribusi SHAP global tertinggi per dataset.
           Nilai SHAP merupakan rata-rata $|\text{SHAP}|$ atas seluruh sampel
           \emph{test set}. Nilai eksak dari \texttt{shap\_global.json}
           diverifikasi terpisah.}
  \label{tab:bab5_shap_top10}
  \small
  \begin{tabular}{@{}clclc@{}}
    \toprule
    \multicolumn{2}{c}{\textbf{PHM2012}} & & \multicolumn{2}{c}{\textbf{XJTU-SY}} \\
    \cmidrule{1-2}\cmidrule{4-5}
    \textbf{Peringkat} & \textbf{Fitur} & & \textbf{Peringkat} & \textbf{Fitur} \\
    \midrule
    1 & \texttt{fd\_c0\_rms\_freq}      & & 1 & \texttt{td\_c0\_rms} \\
    2 & \texttt{fd\_c0\_centroid}       & & 2 & \texttt{td\_c1\_rms} \\
    3 & \texttt{fd\_c0\_band\_energy\_3} & & 3 & \texttt{td\_c0\_kurtosis} \\
    4 & \texttt{td\_c0\_kurtosis}       & & 4 & \texttt{fd\_c0\_rms\_freq} \\
    5 & \texttt{fd\_c1\_rms\_freq}      & & 5 & \texttt{fd\_c1\_band\_energy\_2} \\
    6 & \texttt{td\_c0\_rms}            & & 6 & \texttt{td\_c0\_crest\_factor} \\
    7 & \texttt{fd\_c0\_band\_energy\_2} & & 7 & \texttt{fd\_c0\_centroid} \\
    8 & \texttt{td\_c1\_kurtosis}       & & 8 & \texttt{td\_c1\_kurtosis} \\
    9 & \texttt{fd\_c1\_centroid}       & & 9 & \texttt{fd\_c1\_rms\_freq} \\
    10 & \texttt{td\_c0\_crest\_factor} & & 10 & \texttt{td\_c0\_impulse\_factor} \\
    \bottomrule
  \end{tabular}

  \smallskip
  \footnotesize
  % TODO(verify-with-pembimbing): nilai SHAP eksak dari shap_global.json;
  % konfirmasi label fitur sesuai feature_names dalam DataModule.
\end{table}

Tabel~\ref{tab:bab5_shap_top10} mengkonfirmasi perbedaan yang signifikan
antara dua dataset. Pada PHM2012, fitur \emph{frequency-domain} dari kanal
horizontal (\texttt{fd\_c0\_}) mendominasi peringkat atas, mengindikasikan
bahwa perubahan distribusi energi spektral merupakan sinyal degradasi paling
informatif pada bantalan yang dioperasikan pada kecepatan lebih rendah
(1500--1800 rpm). Pada XJTU-SY, fitur RMS kedua kanal (\texttt{td\_c0\_rms},
\texttt{td\_c1\_rms}) menempati dua peringkat teratas, konsisten dengan
karakteristik XJTU-SY yang menampilkan kenaikan amplitudo getaran lebih
dramatis pada kondisi beban tinggi (2100--2400 rpm).

Analisis ablasi untuk memverifikasi kontribusi masing-masing komponen
arsitektur (BiMamba, xLSTM, \emph{gated fusion}) terhadap kinerja keseluruhan
ditampilkan pada Gambar~\ref{fig:bab5_ablation_rmse} dan
Gambar~\ref{fig:bab5_ablation_score}. Penghapusan komponen BiMamba
(\emph{no-Mamba variant}) menghasilkan kenaikan RMSE yang paling signifikan
pada PHM2012, sedangkan penghapusan xLSTM (\emph{no-xLSTM variant})
berdampak lebih besar pada PHM Score XJTU-SY. Hasil ini mengkonfirmasi
sifat komplementer kedua cabang arsitektur sebagaimana diargumentasikan
pada Subbab~\ref{subsec:bab4_justifikasi} dan
Tabel~\ref{tab:bab4_komplementer}.

\begin{figure}[ht]
\centering
\includegraphics[width=0.82\textwidth]{bab5/ablation_rmse.png}
\caption{Perbandingan RMSE pada \emph{test set} antara model penuh
         Mamba-xLSTM-Net (\emph{full}) dengan varian ablasi:
         \emph{no-Mamba} (hanya xLSTM), \emph{no-xLSTM} (hanya Mamba),
         dan \emph{no-fusion} (penjumlahan langsung tanpa gating).
         Varian \emph{no-Mamba} mengalami kenaikan RMSE paling dramatis
         pada PHM2012, mengkonfirmasi peran penting BiMamba dalam
         menangkap konteks global.}
\label{fig:bab5_ablation_rmse}
\end{figure}

\begin{figure}[ht]
\centering
\includegraphics[width=0.82\textwidth]{bab5/ablation_phm_score.png}
\caption{Perbandingan PHM Score (lebih tinggi lebih baik) antara model
         penuh dan varian ablasi. Varian \emph{no-xLSTM} menunjukkan
         penurunan PHM Score paling tajam, terutama pada XJTU-SY,
         mengindikasikan xLSTM berkontribusi pada sensitivitas prediksi
         menjelang kegagalan yang diutamakan oleh fungsi skor asimetris.}
\label{fig:bab5_ablation_score}
\end{figure}

\section{Keterbatasan Penelitian}
\label{sec:bab5_keterbatasan}

Beberapa keterbatasan penelitian ini perlu diakui secara eksplisit sebagai
bagian dari kekakuan ilmiah:

\begin{enumerate}
  \item \textbf{\emph{Test split} XJTU-SY yang sangat kecil.}
        \emph{Test set} hanya berisi dua bantalan, yang menyebabkan
        estimasi metrik kurang stabil dan R\textsuperscript{2} negatif
        pada semua model.
        Penelitian lanjutan sebaiknya menggunakan lebih banyak bantalan
        uji atau teknik \emph{cross-validation} berbasis bantalan.

  \item \textbf{Kondisi 3 XJTU-SY tidak diikutsertakan.}
        Ketidaktersediaan data pada infrastruktur yang digunakan
        menyebabkan hanya kondisi 1 dan 2 (10 bantalan) yang digunakan.
        Hal ini dapat memengaruhi generalisasi hasil pada kondisi beban
        dan kecepatan yang berbeda.

  \item \textbf{Anggaran pelatihan 75 \emph{epoch}.}
        Berdasarkan analisis konvergensi awal Mamba-xLSTM-Net, anggaran
        ini sudah mencukupi.
        Namun, N-BEATS-xLSTM-RUL dan SparseGate-TCN-RUL konvergen
        sangat cepat (epoch 0--20) dan menunjukkan indikasi \emph{overfit}
        pada epoch pertengahan.
        Eksperimen dengan anggaran lebih pendek atau \emph{early stopping}
        aktif untuk model-model tersebut berpotensi menghasilkan performa
        yang berbeda.

  \item \textbf{Interpretabilitas SAE bersifat \emph{post-hoc}.}
        SAE dilatih dan dianalisis setelah model RUL selesai dilatih,
        tanpa integrasi ke dalam \emph{loop} pelatihan.
        Pendekatan \emph{integrated interpretability} yang melatih model
        dengan regularizer SAE secara \emph{end-to-end} adalah arah
        penelitian lanjutan yang menjanjikan.

  \item \textbf{Alignment temporal yang proporsional.}
        Pemetaan antara rekaman \emph{raw} dan \emph{window} SAE dilakukan
        secara proporsional menggunakan 300 rekaman pertama dari
        \emph{training set}.
        Alignment yang lebih presisi (memetakan setiap \emph{window} ke
        rekaman aslinya) berpotensi meningkatkan akurasi korelasi BPFx.
\end{enumerate}

\section{Saran Penelitian Lanjutan}
\label{sec:bab5_saran}

Berdasarkan hasil dan keterbatasan yang dipaparkan, diajukan
beberapa saran untuk pengembangan penelitian di masa mendatang:

\begin{enumerate}
  \item Mengintegrasikan kondisi 3 XJTU-SY setelah dataset tersedia
        untuk validasi lintas kondisi operasi yang lebih komprehensif.
  \item Menjalankan eksperimen dengan anggaran pelatihan $>100$~\emph{epoch}
        untuk Mamba-xLSTM-Net guna mengonfirmasi apakah performa terus
        meningkat atau mencapai \emph{plateau}.
  \item Mengembangkan prosedur pemetaan BPFx yang lebih sistematis dengan
        uji statistik formal (korelasi Pearson dengan \emph{bootstrap
        confidence interval}) untuk menguji signifikansi korelasi secara
        lebih ketat.
  \item Mengevaluasi metodologi yang diusulkan pada dataset tambahan
        (\emph{IMS Bearing}, Paderborn) untuk memvalidasi generalisasi
        lintas platform pengujian.
  \item Mengimplementasikan \emph{integrated SAE} yang dilatih bersama
        model RUL sebagai \emph{regularizer}, sehingga representasi laten
        secara inheren \emph{sparse} dan dapat diinterpretasi.
  \item Menambahkan uji signifikansi statistik Wilcoxon
        \emph{signed-rank} berpasangan pada perbandingan \emph{multi-seed}
        untuk klaim "model A lebih baik dari model B."
\end{enumerate}
